\section{Estratégias de projeto}
\label{sec:estrategias}

\status{Pronta. Todas as decisões descritas estão implementadas; as que foram
tomadas após medida trazem a medida. A seção \ref{sec:est-escala} traz números de
escalabilidade em duas dimensões e a seção \ref{sec:par-escala} os de três,
todos medidos nesta máquina.}

\subsection{A floresta é a estrutura}

A decisão da qual as outras dependem. No t8Flow a malha \textbf{é} a floresta do
t8code: a topologia, a numeração, os vizinhos e a partição vêm dela, e nenhuma
estrutura paralela as duplica. Isso tem três consequências imediatas:

\begin{itemize}
\item \textbf{nenhum processo detém a malha inteira.} A curva de preenchimento
      de espaço (ordem de Morton) particiona as folhas, e cada processo conhece
      as suas mais uma camada de fantasmas de face;
\item \textbf{o refino é da floresta.} Adaptar é chamar o adaptador do t8code e
      \emph{reconfirmar} o balanceamento 2:1; não há sincronização entre duas
      estruturas que possam divergir;
\item \textbf{a partição é do t8code.} Não há cálculo próprio de balanceamento de
      carga. Essa é uma renúncia deliberada, e a seção \ref{sec:est-escala} a
      quantifica.
\end{itemize}

\subsection{Multiblocos é multiárvore}

O t8Flow é multiblocos, como o antecessor, e essa foi a peça que encaixou sem
atrito: a malha grossa (\emph{cmesh}) do t8code é nativamente multiárvore, e
\textbf{um bloco é uma árvore}. Blocos vizinhos são unidos pelas faces
coincidentes; uma sobreposição parcial de faces é recusada com exceção, porque
unir faces que não coincidem exatamente produziria uma topologia plausível e
errada.

\subsection{Um domínio só}
\label{sec:est-dominio}

Esta decisão vem de um defeito concreto do antecessor, e por isso é relatada com
ele.

No antecessor, o domínio da fração volumétrica do método VOF é, por construção,
uma malha \emph{uniforme no nível mais fino}, separada da árvore adaptativa que
carrega a quantidade de movimento. Sete funções de advecção tridimensional
percorriam a árvore adaptativa e usavam o índice obtido ali \emph{diretamente}
como índice de uma propriedade que vive no outro domínio. As contagens medidas
nos campos salvos:

\begin{center}
\begin{tabular}{lr}
\toprule
propriedade & células \\
\midrule
fração volumétrica (domínio uniforme) & \num{375000} \\
velocidade (árvore adaptativa)        & $\approx$ \num{81000} \\
\bottomrule
\end{tabular}
\end{center}

\textbf{Em malha uniforme os dois domínios coincidem célula a célula, e o índice
errado é acidentalmente o certo.} Foi por isso que todo o conjunto de testes
tridimensionais passou por anos. Com um nível de refino, oitenta mil iterações
escreviam em um vetor de trezentos e setenta e cinco mil posições sem relação
geométrica com a célula de origem --- e a fração parou de advectar, em silêncio,
enquanto a quantidade de movimento continuava correta.

A decisão para o sistema novo é, portanto, \textbf{um domínio só}: toda
propriedade vive na mesma floresta. Dois domínios paralelos exigiriam tradução de
índice explícita e verificada, e herdar o arranjo sem a tradução herdaria a
classe inteira de defeito. Paga-se barato no início e caro depois.

\subsection{Herança do que é independente de malha}

Duas peças do antecessor~\cite{sousa2019} foram reaproveitadas sem modificação, depois de
verificado que não dependem da estrutura de malha:

\begin{description}
\item[o invólucro de solvers lineares,] com implementações para PETSc,
      \emph{hypre}, SOR, um solver centralizado e um que escreve o sistema em
      disco. A interface recebe linhas, um lado direito, e devolve a solução;
      nada nela conhece malha;
\item[o aproximador por mínimos quadrados móveis ponderados] (seção
      \ref{sec:aprox-mls}) --- a interpolação sem malha que é a contribuição
      central de~\cite{sousa2019} e o que permite àquele método operar em
      árvores \emph{não graduadas} ---, que recebe uma nuvem de pontos e um
      ponto-alvo e devolve pesos. Também não conhece malha.
\end{description}

Uma verificação metodológica vale registro, porque a primeira tentativa dela foi
\emph{inválida}: para decidir se um arquivo é independente de malha, tentou-se
compilá-lo com um diretório de cabeçalhos-sombra vazios. O teste não provou nada,
porque a diretiva \verb|#include "x.h"| procura primeiro o diretório do próprio
arquivo-fonte --- de modo que a sombra nunca se aplicava, e um arquivo que
\emph{depende} de malha foi declarado independente. O conserto foi copiar cada
fonte para fora da árvore antes de compilar.

\subsection{Duas e três dimensões em bases separadas}

O t8code trata malhas bidimensionais e tridimensionais; o t8Flow mantém
\textbf{bases de código separadas} para cada uma, começando pela
bidimensional. A alternativa --- uma base só, com a dimensão como parâmetro de
compilação --- economiza linhas e cobra em outro lugar: cada laço sobre direções,
cada contagem de vizinhos de face e cada fórmula de quadratura ganha um caso que
não é exercitado no regime em que se está trabalhando. Com bases separadas, o
código bidimensional é lido e testado como código bidimensional.

\subsection{Armazenamento: centros e facetas}

As propriedades escalares e tensoriais vivem nos \textbf{centros} das células; as
vetoriais, nas \textbf{facetas} --- o arranjo deslocado (MAC), que é o que torna
o acoplamento pressão--velocidade estável sem estabilização adicional.

A numeração de facetas exige uma decisão de propriedade, e ela é: a faceta
pertence à célula para a qual é a face \emph{``menos''} em cada direção, e
exatamente a uma. As faces ``mais'' que caem no contorno do domínio recebem
faceta própria, porque não há célula à direita que as possua. Em uma malha
uniforme de $N\times N$ isso dá $N(N+1)$ facetas por direção, e essa contagem é
afirmada pela suíte --- é barata e pega erro de propriedade antes que ele apareça
como divergência que não fecha.

\subsection{Estênceis em lote, e a GPU quando houver}
\label{sec:est-gpu}

O estêncil de ordem arbitrária é pedido ponto a ponto ou \textbf{em lote}: $N$
pontos de consulta, $N$ estênceis, em uma chamada. A forma de pedir foi definida
antes do \emph{back-end}, porque é ela que decide se a execução em acelerador é
possível.

\begin{description}
\item[seleção do suporte] fica sempre na CPU --- é navegação de malha, busca por
      vizinhança, crescimento por anéis;
\item[a resolução dos sistemas normais] é o que vai para o acelerador, com
      implementação em \textbf{OpenCL} (escolhido por portabilidade: funciona com
      placas AMD via \emph{rusticl}, Intel via NEO, e NVIDIA);
\item[o \emph{back-end} de CPU existe sempre,] e a escolha entre os dois é
      \textbf{sinalizador de compilação}, não detecção silenciosa em tempo de
      execução.
\end{description}

Precisão simples basta para o estêncil, e isso foi uma decisão consciente: os
pesos entram depois em um sistema linear resolvido em precisão dupla, e o erro
dos pesos é de truncamento do ajuste, não de arredondamento.

O oráculo desse caminho é \emph{ponto a ponto contra a CPU}, e ele encontrou um
defeito que a reprodução polinomial não encontrava: o núcleo em OpenCL usava
mínimos quadrados \textbf{não ponderados}, enquanto o núcleo herdado usa o peso
móvel $1/(d^2+\varepsilon)$. Ambos reproduzem polinômios do grau da ordem
exatamente --- e por isso o teste de reprodução passava nos dois ---, mas com
\emph{pesos diferentes}: diferença medida de \num{0.311} em um peso individual.
Ponderado corretamente, a diferença caiu para \num{8.8e-11}.

Dois percalços de medição no mesmo caminho, registrados porque são o tipo de erro
que faz um relatório afirmar o contrário do que mediu. Primeiro, recompilar com o
sinalizador trocado \emph{não} reconstruía, porque o sinalizador não estava nas
dependências do \emph{make}; o conserto foi gravar a configuração em um arquivo
conferido na leitura do \emph{makefile}. Segundo, a diferença GPU--CPU medida deu
\textbf{exatamente zero} --- porque com a GPU compilada ela é o \emph{back-end}
padrão, e a ``referência de CPU'' estava sendo calculada na GPU. O caso agora
força a CPU explicitamente e \emph{rejeita} o zero exato como suspeito.

\subsection{A escalabilidade, medida}
\label{sec:est-escala}

Duas medidas, nesta máquina, em duas dimensões.

A primeira é o \textbf{ganho de divisão}: \num{9377} células em quatro processos
renderam apenas \num{1.6}$\times$. A leitura não é que não escala --- é que
\num{2300} células por processo é pouco trabalho para o custo de comunicação. O
regime em que a divisão vence é o de mais células por processo, e ele é limitado
por \emph{tamanho}, não por dimensão.

O regime tridimensional, onde o antecessor tem \num{375000} células --- quarenta
vezes mais trabalho por processo ---, passou a ser medido e está na seção
\ref{sec:par-escala}: \num{4.92}$\times$ em oito processos no esquema de fluxo,
com \num{262144} células e erro idêntico ao serial.

\reserva{essa medida foi feita em \textbf{oito núcleos físicos}, e é o teto
desta máquina. O escalonamento além disso --- que é onde o regime do antecessor
de fato vive --- continua sendo expectativa, não resultado.}

A segunda é o \textbf{desbalanceamento por refino local}: refinar aleatoriamente
desbalanceia a partição em \num{2.8}$\times$, e reparticionar corrige isso a um
custo de 3 a 6\% do tempo de passo.

\subsection{A suíte nasce com o código}

A última estratégia é de processo, e foi pedida explicitamente: a suíte de testes
nasce junto com o código, e não depois dele. A forma que isso toma aqui é um
\textbf{contrato} --- um cabeçalho com cláusulas numeradas, cada uma declarando
uma propriedade que o sistema afirma e nomeando o caso que a verifica. O arreio
de execução \emph{lê as cláusulas do contrato} e exige um caso passando para cada
uma: cláusula sem teste sai como falha, não como omissão. A seção
\ref{sec:verificacao} descreve o estado e o que ele já pegou.
